\documentclass[journal,10pt]{IEEEtran}
\setlength{\textwidth}{182mm}
\setlength{\columnsep}{4.2mm}
\usepackage{amsmath,amssymb,newtxmath}
\usepackage[no-math]{fontspec}
\setmainfont{Times New Roman}
\usepackage{xeCJK}
\setCJKmainfont[AutoFakeBold=2]{SimSun}
\usepackage{graphicx,booktabs,array,hyperref,placeins}
\hypersetup{hidelinks,unicode=true}
\raggedbottom
\setlength{\emergencystretch}{2em}
\renewcommand{\arraystretch}{1.15}
\usepackage{seqsplit}
\begin{document}
\title{Litz 线分层建模与数值验证：完整研究报告 · D 数据来源与环境交接}
\author{缪雨轩\thanks{淮南师范学院}}

\maketitle
\noindent 合并稿对应章节 / Chapter in merged manuscript: \href{report_zh-CN_ieee.pdf#chapter-D}{D}\par
\FloatBarrier
\section*{D 数据来源与环境交接}
\hypertarget{chapter-D}{}
\pdfbookmark[0]{D 数据来源与环境交接}{bookmark-chapter-D}
复现分为研究计算与成稿构建两个入口。研究入口读取原始输入和求解配置，执行相应验证；成稿入口读取冻结数据、已验收单图与 Markdown 正文，生成章节、合并稿和多种格式。每一阶段保存输入哈希和输出哈希，只有依赖与产物一致时才复用已完成阶段。独立单图保留矢量曲线及其字体，组图清单保存面板身份与布局参数。


来源清单 / Source manifest: source\_provenance.json; quantities.json; asset\_sources.json.


\begin{table*}[!htbp]
\centering{\footnotesize 冻结来源与完整性 / Frozen sources and integrity}\par
\begingroup\footnotesize\begin{tabular}{p{84.9mm}|p{84.9mm}}
\toprule
Source & SHA256\\
docs/总控/当前运行索引.json & \seqsplit{3324976d4af8f58e60820a88ae978246443f78fbdd0cfa1050e86809110ae60d}\\
outputs/prewrite\_\allowbreak closure\_\allowbreak state.json & \seqsplit{80a66afa79c321a89f730b102af0e70125f6e1655372717d6b59f470c6520e71}\\
\bottomrule\end{tabular}\endgroup
\end{table*}
\end{document}
